problem:  
Let \((\mathcal{L}^{n+1}, \bar g)\) be a Lorentzian manifold admitting a **timelike conformal Killing field** \(\bar T\),  
\[
\mathcal L_{\bar T} \bar g = 2\psi\, \bar g.
\]
Let \(M^n \subset \mathcal{L}^{n+1}\) be a complete spacelike hypersurface with induced metric \(g\), second fundamental form \(h\), and mean curvature \(H\).  
Assume the **tangential projection** \(T = (\bar T)^T\) satisfies \(T = \nabla f\) for some smooth function \(f\) making \((M,g,f)\) a **gradient shrinking Ricci soliton**:
\[
\mathrm{Ric}_g + \nabla^2 f = \lambda g, \quad \lambda>0.
\]
Denote by \(N\) the unit future normal of \(M\), and write the decomposition \(\bar T = T + \tau N\), where \(\tau:\!M\to \mathbb R\) is the normal component of \(\bar T\) along \(M\).  
The conformal Killing identity induces on \(M\) the fundamental relation
\[
\mathcal L_T g = 2(\psi - \tau H)\, g - 2\tau\, h.
\]

**Conjectural Open Problem:**  
Assume that along \(M\) the conformal factor \(\psi\) and normal component \(\tau\) are functionally dependent, for example \(\psi=\psi(\tau)\).  Determine whether there exist nontrivial configurations \((\mathcal L,\bar g,M,f)\) such that the compatibility of  
\[
\mathcal L_T g = 2(\psi - \tau H) g - 2\tau h
\]
with the soliton equation
\[
\mathrm{Ric}_g + \nabla^2 f = \lambda g
\]
forces one of the following **rigid alternatives**:
1. Either \(M\) is **totally umbilic** (i.e. \(h = \frac{H}{n} g\)) and the induced metric evolves self‑similarly under \(\bar T\);  
2. Or there exists a non‑umbilic family of solutions for which the mixed term \(2\tau h\) produces an intrinsic obstruction equation relating the principal curvatures of \(M\) to the ambient conformal structure—potentially describing genuinely new “Lorentzian‑conformal Ricci soliton” hypersurfaces.

Equivalently, the conjecture asks for a classification or obstruction criterion: **Can a gradient shrinking Ricci soliton arise as a spacelike hypersurface whose intrinsic soliton vector field is the tangential part of a timelike conformal motion, and if so, under which curvature or mean‑curvature conditions?**

---

Why is it a "good" problem:  
This reformulation eliminates ambiguities and pinpoints a genuine geometric tension between two tensorial identities—one intrinsic (\(\mathrm{Ric}_g + \nabla^2 f = \lambda g\)) and one extrinsic (\(\mathcal L_T g = 2(\psi-\tau H) g - 2\tau h\))—linked through the decomposition of a timelike conformal field.  
The heart of the problem is whether these equations can be made **mutually consistent** without collapsing to the trivial umbilic case, that is, whether the soliton structure can survive the nontrivial normal coupling encoded by \(h\).  

It is deeply challenging because it requires blending:
- the **Gauss–Codazzi equations** for spacelike hypersurfaces in Lorentzian manifolds,  
- the **Lie derivative interplay** between the ambient conformal symmetry and the intrinsic metric, and  
- the **analytic structure** of gradient Ricci soliton equations.  

A solution would identify a new rigidity or existence principle connecting Lorentzian conformal geometry and Riemannian soliton theory—completely uncharted territory.  
The formulation is simple to state (a single vector decomposition and two coupled tensor identities) yet profound, asking whether these structures can coexist nontrivially. It thus fits the paradigm of a “narrow entrance, deep depth” problem and provides a clear, logically sound research direction.